

\begin{frame}
  \frametitle{Consultas de caminos: idea del algoritmo}

  \begin{itemize}
    \item Si el camino está completamente contenido en un camino pesado, hacemos la consulta directamente.
    \item Si no, el camino \(u\)--\(v\) necesariamente llega hasta la cima de un camino pesado y se sale hacia arriba.
    \item Detectamos cuál de las dos puntas es la que se sale por la cima, y hacemos la consulta en ese camino pesado.
    \item De esta forma resolvemos la punta del camino \(u\)--\(v\) y achicamos el camino que falta por resolver.
  \end{itemize}

  \vspace{0.75em}

\end{frame}

\begin{frame}[fragile]
  \frametitle{Consultas de caminos: consulta}

\begin{lstlisting}
int query(int u, int v) {
  int tu = htop[u], tv = htop[v];
  if (tu == tv) {
    if (dep[u] > dep[v]) swap(u, v);
    return tree[tu].query(dep[u]-dep[tu], dep[v]-dep[tu]+1);
  } else {
    if (dep[tu] < dep[tv]) swap(u, v), swap(tu, tv);
    return tree[tu].query(0, dep[u]-dep[tu]+1) + query(p[tu], v);
  }
}
\end{lstlisting}
\end{frame}